\subsection{\texorpdfstring{情形 \(\chi = 1\)}{情形 \textbackslash chi = 1}}

若 \(f\) 是定义在 \(\mathbf{X} / \mathbf{H}\) 上的函数，则 \(f \circ \pi\) 是定义在 \(\mathbf{X}\) 上且在轨道上为常值的函数；它连续当且仅当 \(f\) 连续。映射 \(f \mapsto f \circ \pi\) 特别定义了从 \(\mathcal{K}(\mathbf{X} / \mathbf{H})\) 到 \(\mathcal{K}^1(\mathbf{X})\) 的双射。

于是，在 \(\chi = 1\) 的情形，可以将第 1 号的某些结果重新表述如下：

设 \(f\) 是 \(X\) 上的连续数值函数，其支集与 \(X\) 的每个紧子集的饱和集的交都是紧的。公式

\[
f ^ {b} (\pi (x)) = \int_ {\mathrm{H}} f (x \xi) d \beta (\xi)\tag{3}
\]

定义了连续函数 \(f^{b}\) 在 \(\mathbf{X} / \mathbf{H}\) 上。若 \(g\) 是 \(\mathbf{X} / \mathbf{H}\) 上的连续函数，则

\[
(f \cdot g \circ \pi) ^ {b} = f ^ {b} \cdot g.\tag{4}
\]

若 \(\eta \in \mathbf{H}\)，则

\[
\left(\delta (\eta) f\right) ^ {b} = \Delta_ {\mathrm{H}} (\eta) ^ {- 1} f ^ {b}.\tag{5}
\]

不要忘记 \(f^{b}\) 的定义取决于 \(\beta\) 的选择。若 H 是紧的且 \(\beta\) 已归一化，函数 \(f^{b}\) 有时称为 f 的轨道平均。

若 \(f \in \mathcal{K}(\mathrm{X})\)，则 \(f^{b} \in \mathcal{K}(\mathrm{X / H})\)。映射 \(f \mapsto f^{b}\) 从 \(\mathcal{K}(\mathrm{X})\) 到 \(\mathcal{K}(\mathrm{X / H})\) 是线性的，且 \(\mathcal{K}(\mathrm{X})\)（分别为 \(\mathcal{K}_{+}(\mathrm{X})\)）的像是 \(\mathcal{K}(\mathrm{X / H})\)（分别为 \(\mathcal{K}_{+}(\mathrm{X / H})\)）。

评注 1. --- 我们将证明映射 \(f \mapsto f^{b}\) 是从 \(\mathcal{K}(\mathrm{X})\) 到 \(\mathcal{K}(\mathrm{X}/\mathrm{H})\) 的严格态射（GT, III, §2, No.~8）。

\begin{enumerate}
\def\labelenumi{\alph{enumi})}
\item
  该映射连续：只需证明，对于 X 的每个紧子集 K，将 \(f \mapsto f^{b}\) 限制到 \(\mathcal{K}(X, K)\) 后，是从 \(\mathcal{K}(X, K)\) 到 \(\mathcal{K}(X / H, \pi(K))\) 的连续映射（TVS, II, §4, No.~4, 命题 5）；由于 H 在 X 上适当地作用，集合 \(\xi \in H\) 中满足 K \(\xi\) 与 K 相交的元素所成的集合是紧的；由（3）可得 \(\sup_{x \in K} |f^{b}(\pi(x))| \leqslant \beta(P) \sup_{x \in K} |f(x)|\)，这就证明了断言。
\item
  令 \(K'\) 为 X/H 的紧子集。选取 X 中的紧子集 K，使 \(\pi(K) = K'\)，并证明将 \(f \mapsto f^{b}\) 限制到 \(\mathcal{K}(X, K)\) 后，是从 \(\mathcal{K}(X, K)\) 到 \(\mathcal{K}(X/H, K')\) 的严格态射。只需为这个限制构造一个右逆（GT, III, §6, No.~2, 命题 3）。现在，根据第 1 号引理（采用其记号），将以下映射复合即可得到这样的逆：
\end{enumerate}

\(\alpha)\) 映射 \(f' \mapsto f' \circ \pi\) 从 \(\mathcal{K}(\mathrm{X}/\mathrm{H}, \mathrm{K}')\) 到集合 E，其中 E 由 \(\mathcal{K}^1(\mathrm{X})\) 中支集包含于 KH 的函数组成；

\(\beta)\) 从 E 到 E 的映射，把每个 \(g \in \mathbf{E}\) 对应为在 KH 上等于 \(g / u^1\)、在 X - KH 上等于 0 的函数；

\(\gamma)\) 从 E 到 \(\mathcal{K}(\mathbf{X})\) 的映射，把每个函数 \(h \in \mathbf{E}\) 对应为 \(uh\)。

\begin{enumerate}
\def\labelenumi{\alph{enumi})}
\setcounter{enumi}{2}
\tightlist
\item
  这一点成立后，若 \(\mathbf{V}\) 是 \(\mathcal{K}(\mathbf{X})\) 中 0 的凸邻域，则 \(\mathrm{V} \cap \mathcal{K}(\mathrm{X}, \mathrm{K})\) 是 \(\mathcal{K}(\mathrm{X}, \mathrm{K})\) 中 0 的凸邻域，因而 \(\mathrm{V}^{b} \cap \mathcal{K}(\mathrm{X} / \mathrm{H}, \mathrm{K}')\)
\end{enumerate}

是 \(\mathcal{K}(\mathrm{X / H},\mathrm{K}^{\prime})\) 中 0 的凸邻域，根据 \(b)\)，因而 \(\mathbf{V}^b\) 是 \(\mathcal{K}(\mathrm{X / H})\) 中 0 的邻域（TVS, II, §4, No.~4）。证明完毕。

命题 4. --- a) 设 \(\lambda\) 是 X/H 上的测度。存在唯一一个定义在 X 上的测度 \(\lambda^{\sharp}\)，使得

\[
\int_ {\mathrm{X} / \mathrm{H}} f ^ {b} d \lambda = \int_ {\mathrm{X}} f d \lambda^ {\sharp}\tag{6}
\]

对所有 \(f \in \mathcal{K}(\mathrm{X})\) 成立。有 \(\pmb{\delta}(\xi)\lambda^{\sharp} = \Delta_{\mathrm{H}}(\xi)\lambda^{\sharp}\) 对所有 \(\xi \in \mathrm{H}\) 成立。

\begin{enumerate}
\def\labelenumi{\alph{enumi})}
\setcounter{enumi}{1}
\tightlist
\item
  反之，设 \(\mu\) 是 X 上满足 \(\delta(\xi)\mu = \Delta_{\mathrm{H}}(\xi)\mu\) 对所有 \(\xi \in \mathrm{H}\) 成立的测度。存在唯一一个定义在 X/H 上的测度 \(\lambda\)，使得 \(\mu = \lambda^{\sharp}\)。
\end{enumerate}

这是第 1 号命题 3 的特殊情形。

定义 1. --- 在命题 4 的假设和记号下，称 \(\lambda\) 为 \(\mu\) 关于 \(\beta\) 的商，并记为 \(\frac{\mu}{\beta}\) 或 \(\mu/\beta\)。

映射 \(\lambda \mapsto \lambda^{\sharp}\) 从 \(\mathcal{M}(\mathrm{X / H})\) 到 \(\mathcal{M}(\mathrm{X})\)，正是映射 \(f\mapsto f^{b}\) 从 \(\mathcal{K}(\mathrm{X})\) 到 \(\mathcal{K}(\mathrm{X / H})\) 的转置。令 \(\mathfrak{F}\) 为 \(\mathcal{M}(\mathrm{X / H})\) 上的滤子；\(\lim_{\lambda ,\mathfrak{F}}\lambda^{\sharp}(f) = 0\) 对所有 \(f\in \mathcal{K}(\mathrm{X})\) 成立，等价于 \(\lim_{\lambda ,\mathfrak{F}}\lambda (f^{\prime}) = 0\) 对所有 \(f^{\prime}\in \mathcal{K}(\mathrm{X / H})\) 成立；因此，就弱拓扑而言，映射 \(\lambda \mapsto \lambda^{\sharp}\) 是从 \(\mathcal{M}(\mathrm{X / H})\) 到 \(\mathcal{M}(\mathrm{X})\) 的一个线性子空间的同构。这个子空间是弱闭的，因为它由满足 \(\mu \in \mathcal{M}(\mathrm{X})\) 且 \(\delta (\xi)\mu = \Delta_{\mathrm{H}}(\xi)\mu\) 对所有 \(\xi \in \mathrm{H}\) 成立的测度组成。显然，条件 \(\lambda \geqslant 0\) 和 \(\lambda^{\sharp}\geqslant 0\) 等价。

公式（6）类比通常的二重积分记号可写成

\[
\int_ {\mathrm{X}} f (x) d \lambda^ {\sharp} (x) = \int_ {\mathrm{X/H}} d \lambda (\dot {x}) \int_ {\mathrm{H}} f (x \xi) d \beta (\xi) \quad (\dot {x} = \pi (x)).\tag{7}
\]

这涉及记号的滥用：把积分 \(\int_{\mathrm{H}} f(x\xi) d\beta(\xi)\) 看作 \(\dot{x}\) 而不是 x 的函数；只要不会引起混淆，以下将经常采用这种写法。
% semantic-fragments: legacy-input-seam
\relax{}
